\section{D+1到D+5：按龙头战法每天怎么做}

先把话说清楚：下面说的“主板关系股”，是指公司公告、项目名单或同级一手资料能证明它和AI长剧项目有直接关系的主板股。掌阅科技、博纳影业、中广天择当前不属于这类股；它们涨得再好，也只能说明AI内容或影视方向在动，不能直接当作AI长剧龙头。

\subsection*{D+1：先看题材的主板根基}

\textbf{开盘前：}只查两件事：昨晚项目的新公告；主板新出现的关系股。关系股为空时，今天的任务就是看盘，不做“买不到创业板龙头就买主板”的替代动作。

\textbf{盘中：}芒果超媒继续强，只代表题材热度还在。主板关系股第一次涨停时，把它记进观察表；当天不急着把一板当龙头。主板关系股全是一字板时，也不把涨停价当成自己能成交的价格。

\textbf{收盘后：}分三种结果。

\begin{itemize}
  \item 主板关系股仍未出现：AI长剧只算创业板情绪，主板计划空着。
  \item 主板出现一只关系股首板：明天只盯它能不能走成二板。
  \item 主板出现二板，同时又有新的关系股首板：题材开始有梯队，明天才值得比较谁是高度龙头、谁是容量股。
\end{itemize}

\subsection*{D+2：决定要不要把它当主线跟}

\textbf{开盘前：}把D+1的最高板、同方向首板和成交额最大的关系股列成三行。只看这三行，其他传媒股先放一边。

\textbf{盘中：}最高板全天一字、停牌或涨跌停限制时，跳过它，不去找低位影子股。若能正常交易，也不因为开盘冲高就追；盘中只记录价格和成交，真正的判断放到收盘。

\textbf{收盘后：}只有同时看到下面四件事，才把题材升级为“可以继续跟”的状态：

\begin{enumerate}
  \item 有新的官方催化，或者项目有能核验的播出、经营反馈；
  \item 主板最高板继续保持强势，并且新关系股首板持续补位；
  \item 一只主板关系股连续两日成交额靠前，收盘不弱于前收，且全天并非一字；
  \item 关系股至少有三只，当天大多数收涨。
\end{enumerate}

四件事齐了，\textbf{只留一个人选}：能正常交易时，高度龙头优先；高度龙头无法正常交易时，才看成交额靠前的容量股。四件事缺任何一件，继续看，不把题材当主线做。

\subsection*{D+3：看题材是加强，还是只剩一个高标}

\textbf{开盘前：}查第二个官方催化，例如新项目、项目数据或公司明确关系。

\textbf{收盘后：}第二催化来了，最高板还在、容量股还在、新首板也还在，说明题材在加强，可以把前一天留下的唯一人选继续放在计划里。若最高板还涨，其他关系股已经断掉，或者当天未出现新首板，按“高位孤岛”处理：不加新计划，不追补涨。

\subsection*{D+4到D+5：只留还在变强的题材}

到这个阶段，题材必须继续给新东西。连续两个交易日未新增催化、未新增关系股首板，初始题材就算结束。最高板走弱、成交额最大的关系股也掉队时，删除整张AI长剧观察表；不因为股价回落就去赌“低位补涨”。

\begin{exhibitbox}[展品：每天收盘后只填这六项]
\begin{tabularx}{\textwidth}{L{3.0cm}X}
\toprule
要填的内容 & 用来判断什么 \\
\midrule
当天新增的官方消息 & 催化是否延续。 \\
主板关系股数量 & 主板是否真正参与。 \\
最高连板数和新首板数 & 龙头是否有梯队支撑。 \\
主板关系股的成交额排名 & 谁可能承担容量。 \\
收盘是否强于前收 & 强度是否还在。 \\
是否全天一字、停牌或涨跌停 & 理论上的强势能否转化为可交易性。 \\
\bottomrule
\end{tabularx}
\end{exhibitbox}

这套预案的核心在于每天把“看热闹”“可以跟踪”“题材结束”分清楚。主板关系股出现前，主板账户最正确的动作就是等待。
